% !TeX program = xelatex
% 全国大学生数学建模竞赛 2026 冲奖型论文模板（电子提交版）
% 依据：2026 年全国组委会论文格式规范。
% 编译：XeLaTeX（Overleaf 中将 Compiler 设为 XeLaTeX）。
% 提交前：删除所有“写作提示”和占位内容，确认正文不超过 30 页。

\documentclass[UTF8,zihao=5,a4paper,scheme=chinese]{ctexart}

\usepackage[left=2.5cm,right=2.5cm,top=2.5cm,bottom=2.5cm]{geometry}
\usepackage{amsmath,amssymb,amsthm,mathtools,bm}
\usepackage{booktabs,array,tabularx,longtable,multirow}
\usepackage{graphicx,float,caption,subcaption}
\usepackage{enumitem}
\usepackage{xcolor}
\usepackage{listings}
\usepackage[hidelinks]{hyperref}

\hypersetup{
  pdftitle={全国大学生数学建模竞赛论文},
  pdfauthor={},
  pdfsubject={数学建模竞赛论文}
}

\geometry{headheight=14pt,footskip=1.5cm}
\pagestyle{plain}
\setlength{\parindent}{2em}
\setlength{\parskip}{0pt}
\linespread{1.25}
\setlist{nosep,leftmargin=2.8em}
\captionsetup{font=small,labelsep=quad}

% 采用范文常见的中文一级编号与阿拉伯数字下级编号。
\ctexset{
  section={
    name={,、},
    number=\chinese{section},
    format=\centering\heiti\zihao{3},
    beforeskip=1.2ex,
    afterskip=1.0ex
  },
  subsection={
    format=\raggedright\heiti\zihao{4},
    beforeskip=1.0ex,
    afterskip=0.6ex
  },
  subsubsection={
    format=\raggedright\heiti\zihao{-4},
    beforeskip=0.8ex,
    afterskip=0.4ex
  }
}

\definecolor{guidegray}{RGB}{105,105,105}
\definecolor{codegray}{RGB}{247,247,247}
\newcommand{\guide}[1]{%
  \par\noindent{\kaishu\zihao{5}\color{guidegray}【写作提示】#1}\par
}
\newcommand{\keywords}[1]{%
  \par\noindent\textbf{关键词：}#1
}
\newcommand{\placeholder}[1]{%
  \par\noindent{\kaishu\color{guidegray}【#1】}\par
}

\newtheorem{definition}{定义}
\newtheorem{theorem}{定理}
\newtheorem{proposition}{命题}
\newtheorem{assumption}{假设}

\lstdefinestyle{cumcm}{
  basicstyle=\ttfamily\zihao{5},
  backgroundcolor=\color{codegray},
  frame=single,
  rulecolor=\color{black!20},
  numbers=left,
  numberstyle=\tiny\color{guidegray},
  numbersep=8pt,
  breaklines=true,
  showstringspaces=false,
  tabsize=4,
  columns=fullflexible,
  keepspaces=true
}

\begin{document}
\pagenumbering{arabic}
\setcounter{page}{1}

% ==================== 摘要专用页 ====================
\begin{center}
  {\heiti\zihao{2}\bfseries 论文题目（准确揭示研究对象、核心方法与目标）\par}
  \vspace{1.5em}
  {\heiti\zihao{-2}\bfseries 摘\quad 要\par}
\end{center}

\guide{摘要原则上不超过 1 页，完成“任务—方法—结果—检验—价值”的闭环。少写背景，多给可核验的数值结论；模型名称必须具体。}

\noindent\textbf{【总述】}用 2--3 句界定研究对象、核心矛盾与总体技术路线；紧接一句明确全文最重要的创新点或统一建模主线。

\noindent\textbf{【问题一】}写清“建立什么模型—为什么适用—如何求解—得到什么定量结果”，至少给出一个关键数值、误差或比较结论。

\noindent\textbf{【问题二】}说明相对问题一新增的约束、机制或数据，以及方法升级和核心结果，避免只罗列算法名称。

\noindent\textbf{【问题三/四】}说明决策、预测或优化目标，给出最优方案、区间或排序及其现实含义；按实际题目增删。

\noindent\textbf{【收束】}报告验证结果，如误差、拟合、交叉验证、敏感性、扰动试验或基线对比，并说明模型可迁移的边界。

\keywords{模型名称；关键算法；研究对象；创新点（建议 3--5 个）}

\guide{摘要终检：能否脱离正文独立阅读？是否有关键数字？是否说清创新？是否删除学校、队员和赛区信息？}

\clearpage

% ==================== 正文（不设置目录） ====================
\section{问题重述}

\subsection{问题背景}
\guide{用自己的话压缩背景，只保留与变量、目标、约束直接相关的信息，通常 1--2 段。}
\placeholder{填写正文}

\subsection{问题提出}
\guide{逐项重述输入、输出和约束，不照抄原题；明确“已知什么、求什么、评价标准是什么”。}
\begin{enumerate}[label=（\arabic*）]
  \item 问题一：……
  \item 问题二：……
  \item 问题三：……（按赛题实际增删）
\end{enumerate}

\section{问题分析与总体思路}

\subsection{问题间的逻辑关系}
\guide{说明各问是递进、并列还是共享参数，并用技术路线图展示“数据—模型—求解—验证—结论”。}

\begin{figure}[H]
  \centering
  \fbox{\parbox[c][4cm][c]{0.82\textwidth}{\centering 在此插入总体技术路线图}}
  \caption{总体技术路线}
  \label{fig:roadmap}
\end{figure}

\subsection{分问题分析}

\subsubsection{问题一的分析}
\guide{依次写任务本质、关键难点、模型选择依据、预期输出与检验方式，不提前堆公式。}

\subsubsection{问题二的分析}
\guide{说明新增约束、共享参数以及相对上一问的方法升级。}

\subsubsection{问题三的分析}
\guide{说明决策、预测或优化目标，以及结果如何被验证。}

\section{模型假设}
\guide{仅保留不可由题目或数据直接保证、但对模型成立必要的假设，并说明合理性和影响。}

\begin{table}[H]
  \centering
  \caption{模型假设及其依据}
  \label{tab:assumptions}
  \begin{tabularx}{0.94\textwidth}{c X X}
    \toprule
    编号 & 假设内容 & 依据与影响 \\
    \midrule
    1 & …… & …… \\
    2 & …… & …… \\
    3 & …… & …… \\
    \bottomrule
  \end{tabularx}
\end{table}

\section{符号说明}
\guide{仅列跨章节频繁出现的核心符号；局部符号在公式附近定义。}

\begin{table}[H]
  \centering
  \caption{主要符号说明}
  \label{tab:symbols}
  \begin{tabularx}{0.86\textwidth}{c X c}
    \toprule
    符号 & 含义 & 单位/取值域 \\
    \midrule
    $x_i$ & …… & …… \\
    $\theta$ & …… & …… \\
    $f(\cdot)$ & …… & -- \\
    \bottomrule
  \end{tabularx}
\end{table}

\section{数据处理与探索性分析}
\guide{若无外部数据或无需预处理，可将本章合并到各问题章节。}

\subsection{数据来源与可信性}
\placeholder{数据来源、采集日期、字段解释、引用编号与公开性}

\subsection{数据预处理}
\placeholder{缺失值、异常值、重复值、量纲统一、特征构造及处理前后样本量}

\subsection{探索性分析与建模启示}
\placeholder{只保留能支撑后续模型选择的统计量和图，每幅图均说明“看到什么，因此如何建模”}

% 复制下面的“问题模型”结构即可扩展问题数量。
\section{问题一的模型建立与求解}

\subsection{目标、变量与约束}
\guide{一次定义决策变量或状态变量、目标函数、硬约束及评价指标。}

\subsection{模型原理与关键推导}
\guide{从机理、守恒、概率或优化逻辑推导最终模型；所有公式后解释符号与作用。}
\begin{equation}
  \min_{\bm{x}}\; f(\bm{x})
  \quad\text{s.t.}\quad
  g_i(\bm{x})\le 0,\; i=1,2,\ldots,m.
  \label{eq:model}
\end{equation}
式~\eqref{eq:model} 中，$\bm{x}$ 表示……，$f(\bm{x})$ 表示……。

\subsection{求解算法与实现}
\guide{给出算法流程、初始化、终止条件、复杂度和关键参数；创新步骤详述，通用步骤简述。}
\begin{enumerate}[label=步骤 \arabic*：,leftmargin=5em]
  \item 初始化参数和可行解；
  \item 更新状态并计算目标函数；
  \item 判断终止条件，输出最优解及诊断指标。
\end{enumerate}

\subsection{结果呈现与解释}
\guide{先给核心结果表或图，再解释趋势、机制与现实含义；报告精度、单位和有效数字。}
\begin{table}[H]
  \centering
  \caption{问题一核心结果}
  \label{tab:result1}
  \begin{tabularx}{0.94\textwidth}{X c c X}
    \toprule
    指标/方案 & 结果 & 基线/题设要求 & 结论 \\
    \midrule
    …… & …… & …… & …… \\
    …… & …… & …… & …… \\
    \bottomrule
  \end{tabularx}
\end{table}

\subsection{本问检验与小结}
\guide{选择与本问风险匹配的检验，如残差、基线对比、边界案例、消融、交叉验证或扰动实验；用 2--3 句回答本问。}

\section{问题二的模型建立与求解}
\guide{沿用问题一的证据链，但应突出新增约束、模型升级和关键结果，不必机械重复相同小节。}
\subsection{模型改进与新增约束}
\placeholder{填写模型改进、公式推导和参数定义}
\subsection{算法、结果与检验}
\placeholder{填写求解过程、定量结果、基线比较和本问结论}

\section{问题三的模型建立与求解}
\guide{针对决策、预测或优化任务写清目标、约束、方法、结果和验证；按实际问题数增删。}
\subsection{模型建立与求解}
\placeholder{填写模型、算法和实现细节}
\subsection{结果解释与本问检验}
\placeholder{填写最优方案、预测区间或排序结果及稳定性证据}

\section{综合检验与稳健性分析}
\guide{按需保留。若检验已在各问完成且没有跨问参数，可删除本章；若核心参数贯穿全文，集中展示更清晰。}

\subsection{模型有效性与基线对比}
\placeholder{与解析解、已知事实、题设样例、简单基线或替代模型比较}

\subsection{敏感性与鲁棒性}
\placeholder{扰动关键参数、数据或初值，报告输出变化率、置信区间或稳定区域}

\subsection{误差来源与适用边界}
\placeholder{区分数据误差、结构误差与数值误差，并说明模型何时失效}

\section{结论与建议}
\guide{用编号逐一对应题目要求，直接给最终答案；必要时给可执行建议。}
\begin{enumerate}[label=（\arabic*）]
  \item 问题一结论：……
  \item 问题二结论：……
  \item 问题三结论：……
\end{enumerate}

\section{模型评价、改进与推广}

\subsection{模型优势与创新性}
\guide{优势必须有证据，如更低误差、更少参数、更强可解释性、更快计算或更广适用范围。}

\subsection{局限与改进方向}
\guide{写具体限制、影响方向和可执行改进，不写“时间仓促”。}

\subsection{推广与迁移}
\guide{说明可迁移到何种对象、需替换哪些参数或约束、哪些前提必须重新验证。}

% ==================== 参考文献 ====================
\renewcommand{\refname}{参考文献}
\begin{thebibliography}{99}
  \bibitem{ref1} 作者. 题名[J]. 期刊名, 年, 卷(期): 起止页码.
  \bibitem{ref2} 作者. 书名[M]. 版次. 出版地: 出版社, 年.
  \bibitem{ref3} 机构或作者. 网页或数据集名称[EB/OL]. 发布日期[引用日期]. URL.
\end{thebibliography}

% ==================== 附录 ====================
\clearpage
\section*{附录}
\addcontentsline{toc}{section}{附录}

\subsection*{附录 A\quad 支撑材料文件清单}
\guide{清单与最终提交压缩包逐项一致。若无支撑材料，写“本论文没有支撑材料”。}
\begin{table}[H]
  \centering
  \begin{tabularx}{0.96\textwidth}{c X X X}
    \toprule
    序号 & 文件名 & 用途/对应章节 & 运行环境 \\
    \midrule
    1 & …… & …… & …… \\
    2 & …… & …… & …… \\
    \bottomrule
  \end{tabularx}
\end{table}

\subsection*{附录 B\quad 源程序代码}
\guide{附录包含建模所用全部完整、可运行代码。若未使用程序，写“本论文没有用到程序”。}
\begin{lstlisting}[style=cumcm,language=Python,caption={示例代码（提交前替换）}]
def solve_model(data, params):
    """Return model result and diagnostics."""
    result = None
    diagnostics = {}
    return result, diagnostics
\end{lstlisting}

\subsection*{附录 C\quad 补充结果}
\guide{放置不宜占用正文篇幅、但可复核结论的中间结果、完整参数表和补充图。}

\end{document}
